\section{Optimal allocation and minimum-budget laws}
\label{sec:allocation}

\subsection{Budget-constrained design}

Assign each scalable resource component a positive unit cost $c_k>0$ and define
\begin{equation}
 C(\bm e)=\bm c^T\bm e.
\end{equation}
For a total budget $B$, the information-maximization problem is
\begin{equation}
 \mathcal I_{\max}(B;\Lambda)
 =
 \max_{\bm e\ge0,\ \bm c^T\bm e\le B}
 \Ires(\bm e;\Lambda).
\label{eq:max-info}
\end{equation}
The dual resource-closure question is
\begin{equation}
 \Bmin(I_\star;\Lambda)
 =
 \min_{\bm e\ge0}
 \left\{
 \bm c^T\bm e:
 \Ires(\bm e;\Lambda)\ge I_\star
 \right\}.
\label{eq:min-budget}
\end{equation}

\paragraph*{Proposition 2 (convex resource closure).}
Under the assumptions of Proposition 1, Eq.~\eqref{eq:max-info} is a concave maximization over a convex feasible set, and Eq.~\eqref{eq:min-budget} is a convex optimization problem whenever the target superlevel set is nonempty.

\emph{Proof.}
The feasible budget simplex is convex and $\Ires$ is concave. Its superlevel sets
$\{\bm e:\Ires(\bm e;\Lambda)\ge I_\star\}$ are therefore convex. Minimizing the linear cost over such a set is a convex program. \hfill$\square$

For $\Lambda=0$, positive homogeneity makes the budget law especially transparent. Define the best profiled information per unit cost,
\begin{equation}
 \kappa_\star
 =
 \max_{\bm w\ge0,\ \bm c^T\bm w=1}
 \Ires(\bm w;0).
\label{eq:kappa}
\end{equation}
Then
\begin{equation}
 \boxed{
 \mathcal I_{\max}(B;0)=B\kappa_\star,
 \qquad
 \Bmin(I_\star;0)=\frac{I_\star}{\kappa_\star}.
 }
\label{eq:budget-law}
\end{equation}
If $\kappa_\star=0$, no amount of the declared scalable resource can reach a positive local information target. This gives a sharp distinction between a \emph{resource-limited} branch and a \emph{geometry-limited} branch.

\subsection{Marginal information and shadow prices}

Suppose the minimizer $\bm a_\star(\bm e)$ in Eq.~\eqref{eq:resource-variational} is locally unique. The envelope theorem gives
\begin{equation}
 \frac{\partial\Ires}{\partial e_k}
 =
 \|\bm s_k-J_k\bm a_\star(\bm e)\|^2.
\label{eq:marginal}
\end{equation}
Thus the marginal value of another unit of setting $k$ is not its raw target norm $\|\bm s_k\|^2$ but its \emph{residual} target norm after the nuisance coefficient preferred by the current joint design has been fitted.

At a differentiable optimum of Eq.~\eqref{eq:max-info}, Karush--Kuhn--Tucker conditions imply that active settings equalize marginal profiled information per unit cost:
\begin{equation}
 \frac{\|\bm s_k-J_k\bm a_\star\|^2}{c_k}
 =\lambda_B
 \quad\text{for } e_k>0,
\label{eq:kkt-active}
\end{equation}
whereas inactive settings satisfy
\begin{equation}
 \frac{\|\bm s_k-J_k\bm a_\star\|^2}{c_k}
 \le\lambda_B.
\label{eq:kkt-inactive}
\end{equation}
This rule gives an operational meaning to a resource shadow price: a channel is valuable when it contributes a direction that the current nuisance fit cannot cheaply imitate.

\subsection{Analytic continuation of the RQIR-II two-band example}

RQIR II considered a two-band target and relative-tilt nuisance,
\begin{equation}
 \bm g=(g_2,g_4)^T,
 \qquad
 \bm t=(g_2,-g_4)^T.
\end{equation}
Let the two bands now receive independent exposures $e_2,e_4$. Profiling the shared tilt gives
\begin{equation}
 \Ires(e_2,e_4)
 =
 \frac{4A D}{A+D},
 \qquad
 A=e_2g_2^2,\quad D=e_4g_4^2.
\label{eq:weighted-two-band}
\end{equation}
For costs $c_2,c_4$, write cost allocations
$r_2=c_2e_2$, $r_4=c_4e_4$, with $r_2+r_4=B$, and define
\begin{equation}
 \alpha=\frac{g_2^2}{c_2},
 \qquad
 \gamma=\frac{g_4^2}{c_4}.
\end{equation}
The exact optimum is
\begin{align}
 \frac{r_2^\star}{B}
 &=\frac{\sqrt{\gamma}}{\sqrt{\alpha}+\sqrt{\gamma}},
 &
 \frac{r_4^\star}{B}
 &=\frac{\sqrt{\alpha}}{\sqrt{\alpha}+\sqrt{\gamma}},
\label{eq:cost-fractions}\\
 \kappa_\star
 &=\frac{4\alpha\gamma}
 {(\sqrt{\alpha}+\sqrt{\gamma})^2}.
\label{eq:two-band-kappa}
\end{align}
The stronger raw-response band therefore does not necessarily receive the larger resource share. The weaker band can be the expensive-to-replace direction that enables self-calibration.

For the frozen RQIR-II values
\begin{equation}
 g_2=0.3,\qquad g_4=0.7,
\end{equation}
and equal unit costs, a total exposure $B=2$ gives
\begin{equation}
 (e_2^\star,e_4^\star)=(1.4,0.6)
\end{equation}
and
\begin{equation}
 \Ires^\star=0.3528.
\end{equation}
Equal allocation $(1,1)$ gives
\begin{equation}
 \Ires^{\rm equal}=0.3041379310344828.
\end{equation}
Hence optimal nuisance-aware allocation yields
\begin{equation}
 \frac{\Ires^\star}{\Ires^{\rm equal}}-1
 =0.16,
\end{equation}
a $16\%$ information gain at exactly the same total exposure.

As a synthetic budget inversion, if the local design target is $\sigma_\star=0.2$, then $I_\star=25$. Equation~\eqref{eq:budget-law} with
$\kappa_\star=0.1764$ gives
\begin{equation}
 \Bmin=141.7233560090703
\end{equation}
in the same abstract exposure units. This number is deliberately not called seconds, shots, atoms, photons, or joules. Such a label becomes legitimate only after an apparatus-specific resource binding is supplied.

\subsection{Robust allocation under response uncertainty}

A physical resource map may depend on an uncertain but declared response state $\xi\in\Xi$---for example a bounded spectral-density envelope, uncertain visibility, or a calibration-transfer model. Suppose that for every fixed $\xi$, $\Ires(\bm e;\xi)$ has the concavity property of Proposition 1 on a common resource domain. Define the worst-case information
\begin{equation}
 \Ires^{\rm rob}(\bm e)
 =\inf_{\xi\in\Xi}\Ires(\bm e;\xi).
\label{eq:robust-info}
\end{equation}
Because the hypograph of Eq.~\eqref{eq:robust-info} is the intersection of the hypographs for all $\xi$, $\Ires^{\rm rob}$ is also concave. Thus a worst-case resource-closure constraint
\begin{equation}
 \Ires^{\rm rob}(\bm e)\ge I_\star
\end{equation}
remains convex under the stated common-domain assumptions. An average over a declared distribution of $\xi$ is likewise concave. This allows robustness to enter the resource layer without being confused with an unpriced favorable calibration assumption.
